% subject: physics
% type: subjective
% has_diagram: False

% subject: physics
% type: mcq_sc
% has_diagram: False

\item An electron (mass $9 \times 10^{-31}\,\mathrm{kg}$ and charge $1.6 \times 10^{-19}\,\mathrm{C}$) moving with speed $c/100$ ($c$ = speed of light) is injected into a magnetic field $\vec{B}$ of magnitude $9 \times 10^{-4}\,\mathrm{T}$ perpendicular to its direction of motion. We wish to apply an uniform electric field $\vec{E}$ together with the magnetic field so that the electron does not deflect from its path. Then (speed of light $c = 3 \times 10^8\,\mathrm{m s}^{-1}$)
\begin{tasks}(1)
    \task $\vec{E}$ is parallel to $\vec{B}$ and its magnitude is $27 \times 10^2\,\mathrm{V m}^{-1}$
    \task $\vec{E}$ is parallel to $\vec{B}$ and its magnitude is $27 \times 10^4\,\mathrm{V m}^{-1}$
    \task $\vec{E}$ is perpendicular to $\vec{B}$ and its magnitude is $27 \times 10^4\,\mathrm{V m}^{-1}$
    \task $\vec{E}$ is perpendicular to $\vec{B}$ and its magnitude is $27 \times 10^2\,\mathrm{V m}^{-1}$ \ans
\end{tasks}
\begin{solution}
\begin{align*}
v &= \frac{c}{100} \\
  &= \frac{3 \times 10^8}{100} \\
  &= 3 \times 10^6 \ \mathrm{m s}^{-1} \\
\intertext{For no deflection, net Lorentz force must be zero. Hence}
q\vec{E} + q\left( \vec{v} \times \vec{B} \right) &= 0 \\
\vec{E} &= -\left( \vec{v} \times \vec{B} \right) \\
\intertext{Thus $\vec{E}$ must be along the direction opposite to $\vec{v} \times \vec{B}$, so it is perpendicular to $\vec{B}$. Its magnitude is}
E &= v B \\
  &= \left( 3 \times 10^6 \right)\left( 9 \times 10^{-4} \right) \\
  &= 27 \times 10^2 \ \mathrm{V m}^{-1} \\
\intertext{So $\vec{E}$ is perpendicular to $\vec{B}$ and has magnitude $27 \times 10^2\,\mathrm{V m}^{-1}$. Therefore, the correct option is (d).}
\end{align*}
\end{solution}

\begin{idea}
    \begin{align*}
    \intertext{\textbf{Concept:} No deflection requires the net Lorentz force to vanish.}
    \vec{F} &= q\left(\vec{E}+\vec{v}\times\vec{B}\right) \\
    \vec{F} &= \vec{0} \\
    \vec{E} &= -\,\vec{v}\times\vec{B} \\
    \intertext{\textbf{Application:} For $\vec{v}\perp\vec{B}$, the electric field must oppose the magnetic part of the force.}
    E &= \left|\vec{v}\times\vec{B}\right| \\
    &= vB \\
    \intertext{\textbf{Technique:} Set $\vec{E}$ perpendicular to $\vec{B}$ (and also to $\vec{v}$) so that electric and magnetic forces exactly cancel.}
    \end{align*}
\end{idea}
